\chapter{Secuencias automáticas, combinatoria y ecuaciones funcionales}

\section{Números Odiosos y Malvados (Evil numbers)}

Hasta ahora, la Zeta-Regularización se ha aplicado a secuencias con una estructura aritmética clásica: progresiones aritméticas (números naturales, enteros desplazados), secuencias multiplicativas (números primos) o espectros de operadores diferenciales. Sin embargo, el poder de la regularización analítica se extiende mucho más allá de la teoría de números tradicional, adentrándose en la combinatoria y la teoría de la información. En esta sección, introducimos dos conjuntos de enteros definidos no por sus propiedades de divisibilidad, sino por la estructura de su representación en base 2: los números odiosos y los números malvados (evil numbers). Estos conjuntos son el puente natural hacia el estudio de las \textit{series de Dirichlet automáticas}.

\subsection{Definición aritmética y partición de los enteros}

Sea $s_2(n)$ la función que cuenta la suma de los dígitos de la expansión binaria de un entero no negativo $n$. Esta función es ampliamente conocida en teoría de números y combinatoria como el \textit{peso de Hamming} de $n$. Basándonos en la paridad de $s_2(n)$, clasificamos los enteros no negativos en dos conjuntos disjuntos:

\begin{itemize}
    \item \textbf{Números Odiosos ($O$):} Son aquellos enteros positivos cuya expansión binaria contiene un número \textit{impar} de unos. Formalmente,
    $$ O = \{n \in \mathbb{N} : s_2(n) \equiv 1 \pmod 2\} $$
    Los primeros elementos de este conjunto son: $1, 2, 4, 7, 8, 11, 13, 14, 16, 19, 21, 22, 25, 26, 28, 31, \dots$
    
    \item \textbf{Números Malvados ($E$):} Son aquellos enteros positivos cuya expansión binaria contiene un número \textit{par} de unos. Formalmente,
    $$ E = \{n \in \mathbb{N} : s_2(n) \equiv 0 \pmod 2\} $$
    Los primeros elementos de este conjunto son: $3, 5, 6, 9, 10, 12, 15, 17, 18, 20, 23, 24, 27, 29, 30, \dots$
\end{itemize}

Es evidente que ambos conjuntos forman una partición exacta de los enteros positivos: $O \sqcup E = \mathbb{N}$. A diferencia de los números primos, cuya distribución está gobernada por el Teorema de los Números Primos y posee una densidad asintótica nula, los números odiosos y malvados tienen una densidad asintótica exactamente igual a $1/2$. Es decir, la función de conteo espectral para ambos conjuntos satisface $N_O(x) \sim x/2$ y $N_E(x) \sim x/2$. 

Esta densidad lineal implica que las series de Dirichlet asociadas a estos conjuntos tendrán una abscisa de convergencia $\sigma_c = 1$, al igual que la función Zeta de Riemann. Sin embargo, la ausencia de una estructura multiplicativa (el producto de dos números odiosos no es necesariamente odioso) obliga a desarrollar nuevas herramientas analíticas para estudiar sus funciones zeta.

\subsection{La función característica y la Secuencia de Thue-Morse}

Para analizar analíticamente los conjuntos $O$ y $E$, es fundamental introducir sus funciones características y conectarlas con una de las secuencias automáticas más famosas de las matemáticas: la secuencia de Thue-Morse.

Definimos la función característica de los números odiosos, $a(n)$, como:
$$ a(n) = \begin{cases} 1 & \text{si } n \in O \text{ (suma de dígitos impar)} \\ 0 & \text{si } n \in E \text{ (suma de dígitos par)} \end{cases} $$

A partir de $a(n)$, construimos la sucesión $\varepsilon_n$ mediante la transformación algebraica:
$$ \varepsilon_n = (-1)^{a(n)} = 1 - 2a(n) $$

Esta transformación mapea el conjunto $\{0, 1\}$ al conjunto $\{-1, +1\}$. Analicemos el comportamiento de $\varepsilon_n$:
\begin{enumerate}[a)]
    \item Si $n \in E$, entonces $a(n) = 0$, lo que implica $\varepsilon_n = 1$.
    \item Si $n \in O$, entonces $a(n) = 1$, lo que implica $\varepsilon_n = -1$.
\end{enumerate}

La secuencia resultante $(\varepsilon_n)_{n \ge 0}$ comienza con los valores:
$$ 1, -1, -1, 1, -1, 1, 1, -1, -1, 1, 1, -1, 1, -1, -1, 1, \dots $$

Esta es, exactamente, la célebre \textbf{secuencia de Thue-Morse} sobre el alfabeto $\{-1, +1\}$. La secuencia de Thue-Morse posee una propiedad de autosimilitud o recurrencia fractal profunda. Si denotamos el bloque de $2^k$ términos iniciales como $T_k$, la secuencia se construye mediante la regla de sustitución $T_{k+1} = T_k \overline{T_k}$, donde $\overline{T_k}$ es el bloque complementario (intercambiando $1$ por $-1$ y viceversa). 

Esta propiedad combinatoria es la clave que sustituye a la multiplicatividad en el análisis de estas secuencias. Mientras que la función Zeta de Riemann explota la propiedad multiplicativa de los enteros a través del Producto de Euler, las series de Dirichlet construidas sobre la secuencia de Thue-Morse explotarán su \textit{autosimilitud recursiva} para obtener ecuaciones funcionales.

\subsection{Descomposición de las Series de Dirichlet}

El objetivo final de esta sección es preparar el terreno para calcular el producto Zeta-regularizado de todos los números odiosos y malvados. Para ello, debemos construir sus funciones zeta asociadas y demostrar que su continuación analítica depende enteramente de la secuencia de Thue-Morse.

Definimos las series de Dirichlet para los conjuntos $O$ y $E$ como:
$$ \zeta_O(s) = \sum_{n \in O} \frac{1}{n^s}, \qquad \zeta_E(s) = \sum_{n \in E} \frac{1}{n^s} $$

Dado que $O$ y $E$ particionan $\mathbb{N}$, la suma de ambas series debe recuperar la función Zeta de Riemann:
$$ \zeta_O(s) + \zeta_E(s) = \sum_{n=1}^{\infty} \frac{1}{n^s} = \zeta(s) $$

Por otro lado, consideremos la diferencia entre ambas series. Utilizando la función característica $a(n)$ y la relación $\varepsilon_n = 1 - 2a(n)$, podemos reescribir la diferencia de la siguiente manera:
$$ \zeta_O(s) - \zeta_E(s) = \sum_{n=1}^{\infty} \frac{1 - 2a(n)}{n^s} - \sum_{n=1}^{\infty} \frac{2a(n)}{n^s} \quad \text{(Ajustando los índices)} $$
Más directamente, notando que $a(n) = 1$ para $n \in O$ y $0$ para $n \in E$, y recordando que $\varepsilon_n = -1$ para $n \in O$ y $1$ para $n \in E$, tenemos que $a(n) = \frac{1 - \varepsilon_n}{2}$. Sustituyendo esto en la definición de $\zeta_O(s)$:

$$ \zeta_O(s) = \sum_{n=1}^{\infty} \frac{a(n)}{n^s} = \sum_{n=1}^{\infty} \frac{1 - \varepsilon_n}{2n^s} = \frac{1}{2} \zeta(s) - \frac{1}{2} g(s) $$

Donde hemos definido la \textbf{Serie de Dirichlet de Thue-Morse} como:
$$ g(s) = \sum_{n=1}^{\infty} \frac{\varepsilon_n}{n^s} $$

De manera completamente análoga, para los números malvados obtenemos:
$$ \zeta_E(s) = \sum_{n=1}^{\infty} \frac{1 - a(n)}{n^s} = \frac{1}{2} \zeta(s) + \frac{1}{2} g(s) $$

\subsubsection{El problema de la continuación analítica}

Las ecuaciones anteriores son el corazón analítico del problema. Nos revelan que la continuación analítica de $\zeta_O(s)$ y $\zeta_E(s)$ está condicionada exclusivamente a la continuación analítica de la serie $g(s)$. 

La función $\zeta(s)$ es meromorfa en todo el plano complejo, con un polo simple en $s=1$. Por lo tanto, el comportamiento de $\zeta_O(s)$ y $\zeta_E(s)$ depende de si $g(s)$ puede ser continuada analíticamente. A diferencia de las series de Dirichlet multiplicativas, la serie $g(s)$ no posee un Producto de Euler. Sin embargo, la naturaleza \textit{automática} de la secuencia de Thue-Morse (su regla de sustitución $T_{k+1} = T_k \overline{T_k}$) permite traducir la autosimilitud combinatoria en una \textbf{ecuación funcional infinita} para $g(s)$. 

Esta ecuación funcional, que relaciona $g(s)$ con evaluaciones de la misma función en $s+k$ para $k \in \mathbb{N}$, es la herramienta que permite demostrar que $g(s)$ admite una continuación meromorfa a todo el plano complejo $\mathbb{C}$, y que es holomorfa en $s=0$. Este hecho garantiza que $\zeta_O(s)$ y $\zeta_E(s)$ están bien definidas en el origen, abriendo la puerta a la aplicación de la fórmula maestra $\prod \lambda_i = \exp(-\zeta_\Lambda'(0))$ para calcular los productos infinitos de estos conjuntos exóticos.

\section{Series de Dirichlet para secuencias automáticas}

Hasta este punto, la Zeta-Regularización se ha aplicado fundamentalmente a secuencias definidas por propiedades aritméticas clásicas (progresiones, números primos) o espectros de operadores diferenciales. En todos esos casos, la estructura multiplicativa subyacente permitía el uso de herramientas poderosas como el Producto de Euler. Sin embargo, al adentrarnos en la combinatoria y la teoría de la información, nos encontramos con subconjuntos de los enteros definidos por la estructura de sus dígitos en una base dada. Estos conjuntos carecen de multiplicatividad, lo que obliga a reformular nuestro enfoque analítico. En esta sección, construiremos rigurosamente las series de Dirichlet asociadas a los números odiosos y malvados, y demostraremos cómo su destino analítico queda íntimamente ligado a la función Zeta de Riemann a través de la secuencia de Thue-Morse.

\subsection{Construcción de las funciones zeta $\zeta_O(s)$ y $\zeta_E(s)$}

Recordemos que los enteros positivos se particionan en dos conjuntos disjuntos basados en la paridad de la suma de sus dígitos en base 2 (el peso de Hamming): el conjunto de los números odiosos $O$ (suma impar de unos) y el conjunto de los números malvados $E$ (suma par de unos). Dado que $O \sqcup E = \mathbb{N}$, las series de Dirichlet asociadas a estos conjuntos, definidas inicialmente para $\Re(s) > 1$, deben satisfacer la propiedad de aditividad espectral que establecimos en la Clase 1:

$$ \zeta_O(s) + \zeta_E(s) = \sum_{n \in O} \frac{1}{n^s} + \sum_{n \in E} \frac{1}{n^s} = \sum_{n=1}^{\infty} \frac{1}{n^s} = \zeta(s) $$

Para estudiar $\zeta_O(s)$ y $\zeta_E(s)$ de manera individual, necesitamos herramientas algebraicas que nos permitan filtrar los elementos de cada conjunto. Para ello, introducimos la función característica de los números odiosos, $a(n)$, definida como:

$$ a(n) = \begin{cases} 1 & \text{si } n \in O \\ 0 & \text{si } n \in E \end{cases} $$

A partir de $a(n)$, construimos la sucesión auxiliar $\varepsilon_n$ mediante la transformación lineal $\varepsilon_n = (-1)^{a(n)} = 1 - 2a(n)$. Esta transformación mapea el conjunto $\{0, 1\}$ al conjunto $\{-1, +1\}$, de modo que $\varepsilon_n = -1$ si $n$ es odioso, y $\varepsilon_n = +1$ si $n$ es malvado. La secuencia $(\varepsilon_n)_{n \ge 0}$ es, salvo por el primer término (donde $\varepsilon_0 = 1$), la célebre secuencia de Thue-Morse.

Podemos expresar la función característica $a(n)$ directamente en términos de $\varepsilon_n$ como $a(n) = \frac{1 - \varepsilon_n}{2}$. Sustituyendo esto en la definición de $\zeta_O(s)$, obtenemos:

$$ \zeta_O(s) = \sum_{n=1}^{\infty} \frac{a(n)}{n^s} = \sum_{n=1}^{\infty} \frac{1 - \varepsilon_n}{2n^s} = \frac{1}{2} \sum_{n=1}^{\infty} \frac{1}{n^s} - \frac{1}{2} \sum_{n=1}^{\infty} \frac{\varepsilon_n}{n^s} $$

De manera completamente análoga, para los números malvados, cuya función característica es $1 - a(n) = \frac{1 + \varepsilon_n}{2}$, tenemos:

$$ \zeta_E(s) = \sum_{n=1}^{\infty} \frac{1 - a(n)}{n^s} = \frac{1}{2} \sum_{n=1}^{\infty} \frac{1}{n^s} + \frac{1}{2} \sum_{n=1}^{\infty} \frac{\varepsilon_n}{n^s} $$

\subsection{La función zeta de Thue-Morse $g(s)$}

Las ecuaciones anteriores revelan que el análisis de $\zeta_O(s)$ y $\zeta_E(s)$ se reduce al estudio de dos componentes: la función Zeta de Riemann $\zeta(s)$, y una nueva función que codifica exclusivamente la secuencia de Thue-Morse. Definimos esta última como la \textbf{función zeta de Thue-Morse}:

$$ g(s) = \sum_{n=1}^{\infty} \frac{\varepsilon_n}{n^s} $$

Con esta definición, las relaciones fundamentales que vinculan las funciones zeta de los números odiosos y malvados con la función Zeta de Riemann y la función de Thue-Morse toman una forma compacta y elegante:

\begin{itemize}
    \item $\zeta_O(s) = \frac{1}{2} \zeta(s) - \frac{1}{2} g(s)$
    \item $\zeta_E(s) = \frac{1}{2} \zeta(s) + \frac{1}{2} g(s)$
    \item $\zeta_E(s) - \zeta_O(s) = g(s)$
\end{itemize}

Estas identidades son exactas y válidas en el semiplano de convergencia absoluta $\Re(s) > 1$. Sin embargo, su verdadero poder analítico reside en que proporcionan el puente para la continuación analítica. Dado que $\zeta(s)$ es una función meromorfa en todo el plano complejo $\mathbb{C}$ (con un único polo simple en $s=1$), el comportamiento analítico de $\zeta_O(s)$ y $\zeta_E(s)$ queda dictado \textit{exclusivamente} por la capacidad de continuar analíticamente la función $g(s)$.

\subsection{Relación con la función Zeta de Riemann y el desafío analítico}

La descomposición $\zeta_O(s) = \frac{1}{2}\zeta(s) - \frac{1}{2}g(s)$ es el corazón del problema de regularización para secuencias automáticas. En los casos anteriores (como los números naturales o los primos), la estructura multiplicativa de la secuencia permitía factorizar la función zeta en un Producto de Euler, lo que facilitaba enormemente la continuación analítica mediante técnicas de la teoría analítica de números. 

Los conjuntos $O$ y $E$, al estar definidos por la paridad de la suma de dígitos en base 2, carecen por completo de estructura multiplicativa. Por ejemplo, $2 \in O$ y $3 \in O$, pero su producto $6 \in E$. Esta ruptura de la multiplicatividad significa que $g(s)$ no admite un Producto de Euler. No podemos factorizar $g(s)$ en un producto sobre números primos. 

Por lo tanto, nos enfrentamos a un desafío analítico de una naturaleza completamente distinta:

\begin{enumerate}[a)]
    \item \textbf{La abscisa de convergencia de $g(s)$:} A diferencia de $\zeta(s)$, cuyos coeficientes de Dirichlet son todos $1$, los coeficientes de $g(s)$ son $\varepsilon_n \in \{-1, 1\}$. Las sumas parciales de la secuencia de Thue-Morse no crecen linealmente, sino que están acotadas por potencias fraccionarias (relacionadas con la dimensión fractal de la secuencia). Esto implica que la serie de Dirichlet $g(s)$ converge en una región diferente, pero el punto crítico sigue siendo cómo extenderla a $s=0$.
    \item \textbf{El polo en $s=1$:} Dado que la densidad asintótica de los números odiosos y malvados es exactamente $1/2$, ambas funciones $\zeta_O(s)$ y $\zeta_E(s)$ heredan el polo simple de $\zeta(s)$ en $s=1$ con residuo $1/2$. Esto es consistente con el hecho de que $g(s)$ es holomorfa en $s=1$ (de hecho, $g(1) = 0$ por la cancelación alternante de la serie de Thue-Morse).
    \item \textbf{La evaluación en $s=0$:} Para calcular el producto Zeta-regularizado de los números odiosos mediante la fórmula maestra $\prod_{n \in O} n = \exp(-\zeta_O'(0))$, necesitamos evaluar $\zeta_O'(0)$. Según nuestra descomposición, esto requiere calcular $\zeta'(0)$ (que ya conocemos: $-\frac{1}{2}\log(2\pi)$) y $g'(0)$. 
\end{enumerate}

El problema de calcular el producto de todos los números odiosos se ha transformado, gracias a esta relación con la función Zeta de Riemann, en el problema de evaluar la derivada en el origen de la función zeta de Thue-Morse, $g'(0)$. Dado que $g(s)$ carece de multiplicatividad, su continuación analítica y la evaluación de $g'(0)$ no pueden abordarse con las herramientas clásicas de Euler o Riemann. En su lugar, debemos explotar la única propiedad estructural que posee la secuencia: su \textit{autosimilitud combinatoria} (la regla de sustitución de Thue-Morse). Esta propiedad fractal es la que nos permitirá derivar ecuaciones funcionales infinitas para $g(s)$, las cuales serán la llave para su continuación analítica y la posterior evaluación de los productos exóticos.

\section{Ecuaciones funcionales infinitas}

En las secciones anteriores, establecimos que las series de Dirichlet asociadas a los números odiosos y malvados, $\zeta_O(s)$ y $\zeta_E(s)$, pueden descomponerse en términos de la función Zeta de Riemann $\zeta(s)$ y la función zeta de Thue-Morse $g(s) = \sum_{n=1}^{\infty} \varepsilon_n n^{-s}$. Dado que $\zeta(s)$ es meromorfa en todo el plano complejo, el destino analítico de $\zeta_O(s)$ y $\zeta_E(s)$ depende exclusivamente de nuestra capacidad para continuar analíticamente $g(s)$.

A diferencia de las secuencias multiplicativas como los números naturales o los primos, la secuencia de Thue-Morse carece de estructura multiplicativa, lo que imposibilita el uso del Producto de Euler. Sin embargo, la secuencia de Thue-Morse posee una propiedad combinatoria profunda: es una secuencia \textit{2-automática}. Esto significa que sus términos están gobernados por reglas de recurrencia fractal basadas en la expansión binaria de los índices. En esta sección, demostraremos cómo esta autosimilitud combinatoria se traduce directamente en una \textbf{ecuación funcional infinita} que permite la continuación analítica de $g(s)$ a todo el plano complejo.

\subsection{Derivación de la ecuación funcional}

Para derivar la ecuación funcional, partimos de la definición de la serie de Dirichlet de Thue-Morse, válida inicialmente para $\Re(s) > 1$:

$$ g(s) = \sum_{n=1}^{\infty} \frac{\varepsilon_n}{n^s} $$

La estrategia consiste en aprovechar las propiedades de paridad de la secuencia de Thue-Morse. Recordemos que $\varepsilon_n = (-1)^{s_2(n)}$, donde $s_2(n)$ es la suma de los dígitos binarios de $n$. De esto se deducen dos relaciones de recurrencia fundamentales:
\begin{itemize}
    \item $\varepsilon_{2n} = \varepsilon_n$ (multiplicar por 2 en binario solo añade un cero al final, no cambia la suma de los dígitos).
    \item $\varepsilon_{2n+1} = -\varepsilon_n$ (sumar 1 a un número par cambia el último bit de 0 a 1, aumentando la suma en 1 y cambiando el signo).
\end{itemize}

Utilizamos estas propiedades para dividir la serie $g(s)$ en sus partes de índices pares e impares. Dado que $\varepsilon_0 = 1$ (la suma de dígitos de 0 es 0, que es par), podemos escribir:

$$ g(s) = \sum_{n=1}^{\infty} \frac{\varepsilon_{2n}}{(2n)^s} + \sum_{n=0}^{\infty} \frac{\varepsilon_{2n+1}}{(2n+1)^s} $$

Sustituyendo las relaciones de recurrencia:

$$ g(s) = \frac{1}{2^s} \sum_{n=1}^{\infty} \frac{\varepsilon_n}{n^s} - \sum_{n=0}^{\infty} \frac{\varepsilon_n}{(2n+1)^s} $$

El primer término es exactamente $2^{-s} g(s)$. En el segundo término, extraemos el índice $n=0$, que contribuye con $\varepsilon_0 / 1^s = 1$. Así, obtenemos:

$$ g(s) = 2^{-s} g(s) - 1 - \sum_{n=1}^{\infty} \frac{\varepsilon_n}{(2n+1)^s} $$

Reorganizando para aislar $g(s)$:

$$ (1 - 2^{-s}) g(s) = -1 - \sum_{n=1}^{\infty} \varepsilon_n (2n+1)^{-s} $$

El paso analítico crucial es expandir el término $(2n+1)^{-s}$ para relacionarlo con potencias de $n$. Factorizamos $(2n)^{-s}$:

$$ (2n+1)^{-s} = (2n)^{-s} \left(1 + \frac{1}{2n}\right)^{-s} $$

Dado que $\left| \frac{1}{2n} \right| < 1$ para $n \ge 1$, podemos aplicar el teorema del binomio generalizado:

$$ \left(1 + \frac{1}{2n}\right)^{-s} = \sum_{k=0}^{\infty} \binom{-s}{k} \left(\frac{1}{2n}\right)^k = \sum_{k=0}^{\infty} (-1)^k \binom{s+k-1}{k} \frac{1}{(2n)^k} $$

Sustituyendo esta expansión en la suma sobre $n$:

$$ \sum_{n=1}^{\infty} \varepsilon_n (2n+1)^{-s} = \sum_{n=1}^{\infty} \varepsilon_n \sum_{k=0}^{\infty} (-1)^k \binom{s+k-1}{k} \frac{1}{(2n)^{s+k}} $$

Dado que la serie converge absolutamente para $\Re(s) > 1$, podemos intercambiar el orden de las sumatorias:

$$ \sum_{n=1}^{\infty} \varepsilon_n (2n+1)^{-s} = \sum_{k=0}^{\infty} (-1)^k \binom{s+k-1}{k} \frac{1}{2^{s+k}} \left( \sum_{n=1}^{\infty} \frac{\varepsilon_n}{n^{s+k}} \right) $$

Reconocemos que la suma interna es precisamente $g(s+k)$. Por lo tanto:

$$ (1 - 2^{-s}) g(s) = -1 - \sum_{k=0}^{\infty} (-1)^k \binom{s+k-1}{k} \frac{g(s+k)}{2^{s+k}} $$

Para simplificar, extraemos el término $k=0$ de la sumatoria. Para $k=0$, el coeficiente binomial es $\binom{s-1}{0} = 1$, y el término es $2^{-s} g(s)$. Al restar este término de ambos lados, los $2^{-s} g(s)$ se cancelan mágicamente, dejando:

$$ g(s) = -1 - \sum_{k=1}^{\infty} (-1)^k \binom{s+k-1}{k} \frac{g(s+k)}{2^{s+k}} $$

Absorbiendo el signo negativo en la sumatoria, llegamos a la \textbf{ecuación funcional infinita} para la función zeta de Thue-Morse:

$$ g(s) = -1 + \sum_{k=1}^{\infty} (-1)^{k+1} \binom{s+k-1}{k} \frac{g(s+k)}{2^{s+k}} $$

\subsection{Continuación analítica a todo el plano complejo}

La ecuación funcional derivada no es solo una curiosidad algebraica; es la llave maestra para la continuación analítica. Analicemos su estructura:

\begin{enumerate}[a)]
    \item \textbf{Desplazamiento del argumento:} La función $g(s)$ en el lado izquierdo está evaluada en $s$. En el lado derecho, la función $g$ aparece evaluada en $s+k$ para $k \ge 1$. Esto significa que la ecuación expresa el valor de $g(s)$ en términos de sus valores en regiones con parte real estrictamente mayor.
    \item \textbf{Mecanismo de continuación:} Supongamos que $g(s)$ es analítica en un semiplano derecho $\Re(s) > \sigma_0$. La serie original converge absolutamente en esta región. La ecuación funcional nos permite definir $g(s)$ para $\Re(s) > \sigma_0 - 1$, ya que los términos $g(s+k)$ en el lado derecho caerán en la región $\Re(s+k) > \sigma_0 + k - 1 \ge \sigma_0$, donde la función ya está bien definida y es analítica.
    \item \textbf{Convergencia de la serie infinita:} Para un $s$ fijo, el coeficiente binomial $\binom{s+k-1}{k} = \frac{(s+k-1)(s+k-2)\dots s}{k!}$ crece asintóticamente como $\frac{k^{s-1}}{\Gamma(s)}$ para $k$ grande. Sin embargo, el factor $2^{s+k}$ en el denominador proporciona un decaimiento exponencial $2^{-k}$. Esto garantiza que la serie infinita sobre $k$ converge absoluta y uniformemente en cualquier subconjunto compacto del plano complejo $\mathbb{C}$.
\end{enumerate}

Dado que el lado derecho de la ecuación es una suma uniformemente convergente de funciones analíticas (asumiendo que $g(s+k)$ es analítica), el lado derecho es una función entera. Por el principio de identidad analítica, esto demuestra que $g(s)$ puede ser continuada analíticamente a una \textbf{función entera}, es decir, $g(s)$ es holomorfa en todo el plano complejo $\mathbb{C}$, sin ningún polo ni singularidad.

En consecuencia, las funciones $\zeta_O(s) = \frac{1}{2}\zeta(s) - \frac{1}{2}g(s)$ y $\zeta_E(s) = \frac{1}{2}\zeta(s) + \frac{1}{2}g(s)$ heredan esta propiedad, siendo meromorfas en todo $\mathbb{C}$ con un único polo simple en $s=1$ (proveniente de $\zeta(s)$), y siendo perfectamente regulares (holomorfas) en el origen $s=0$.

\subsection{Evaluación en el origen: $g(0)$ y $g'(0)$}

Para aplicar la fórmula maestra de la Zeta-Regularización a los números odiosos y malvados, necesitamos evaluar $g(s)$ y su derivada en $s=0$. La ecuación funcional infinita nos permite hacer esto de manera exacta y rigurosa.

\subsubsection*{Cálculo de $g(0)$}

Evaluamos la ecuación funcional directamente en $s=0$:

$$ g(0) = -1 + \sum_{k=1}^{\infty} (-1)^{k+1} \binom{k-1}{k} \frac{g(k)}{2^k} $$

El coeficiente binomial $\binom{k-1}{k}$ es igual a $\frac{(k-1)(k-2)\dots(0)}{k!} = 0$ para todo entero $k \ge 1$. Por lo tanto, toda la sumatoria se anula idénticamente, y obtenemos:

$$ g(0) = -1 $$

Este es un resultado fundamental. Significa que la función zeta de Thue-Morse no se anula en el origen, sino que toma el valor entero $-1$.

\subsubsection*{Cálculo de $g'(0)$}

Para encontrar $g'(0)$, derivamos la ecuación funcional con respecto a $s$ y evaluamos en $s=0$. Sea $T_k(s)$ el término general de la sumatoria:

$$ T_k(s) = (-1)^{k+1} \binom{s+k-1}{k} \frac{g(s+k)}{2^{s+k}} $$

Recordemos que $\binom{s+k-1}{k} = \frac{s(s+1)\dots(s+k-1)}{k!}$. Este polinomio en $s$ tiene una raíz en $s=0$. Por lo tanto, $T_k(0) = 0$ para todo $k \ge 1$. Al aplicar la regla del producto para derivar $T_k(s)$, el único término que no se anula al evaluar en $s=0$ es aquel donde la derivada actúa sobre el factor $s$ del coeficiente binomial, mientras que los demás factores se evalúan en $s=0$.

La derivada del coeficiente binomial en $s=0$ es:

$$ \left. \frac{d}{ds} \frac{s(s+1)\dots(s+k-1)}{k!} \right|_{s=0} = \frac{1 \cdot 2 \dots (k-1)}{k!} = \frac{(k-1)!}{k!} = \frac{1}{k} $$

Los demás factores de $T_k(s)$ evaluados en $s=0$ son simplemente $2^{-k} g(k)$. Por lo tanto, la derivada del término general en el origen es:

$$ T_k'(0) = (-1)^{k+1} \frac{1}{k} \frac{g(k)}{2^k} $$

Sumando sobre todos los $k \ge 1$, obtenemos la derivada de la función completa:

$$ g'(0) = \sum_{k=1}^{\infty} (-1)^{k+1} \frac{g(k)}{k 2^k} $$

\subsubsection*{Conexión con un producto infinito convergente}

Aunque la expresión anterior para $g'(0)$ es una serie numéricamente convergente, podemos transformarla en una forma más reveladora expandiendo $g(k)$ usando su definición original $g(k) = \sum_{n=1}^{\infty} \varepsilon_n n^{-k}$:

$$ g'(0) = \sum_{k=1}^{\infty} \frac{(-1)^{k+1}}{k 2^k} \sum_{n=1}^{\infty} \frac{\varepsilon_n}{n^k} $$

Dado que las series convergen absolutamente, intercambiamos el orden de sumación:

$$ g'(0) = \sum_{n=1}^{\infty} \varepsilon_n \sum_{k=1}^{\infty} \frac{(-1)^{k+1}}{k} \left(\frac{1}{2n}\right)^k $$

Reconocemos la suma interna como la expansión en serie de Taylor de la función logaritmo natural, $\log(1+x) = \sum_{k=1}^{\infty} \frac{(-1)^{k+1}}{k} x^k$, evaluada en $x = \frac{1}{2n}$. Por lo tanto:

$$ g'(0) = \sum_{n=1}^{\infty} \varepsilon_n \log\left(1 + \frac{1}{2n}\right) = \sum_{n=1}^{\infty} \varepsilon_n \log\left(\frac{2n+1}{2n}\right) $$

Utilizando las propiedades de los logaritmos, esto se puede escribir como el logaritmo de un producto infinito:

$$ g'(0) = \log \prod_{n=1}^{\infty} \left(\frac{2n+1}{2n}\right)^{\varepsilon_n} $$

Si definimos la constante $Q$ como el producto infinito inverso:

$$ Q = \prod_{n=1}^{\infty} \left(\frac{2n}{2n+1}\right)^{\varepsilon_n} $$

Entonces, hemos demostrado rigurosamente que:

$$ g'(0) = -\log Q $$

\subsection{Significado analítico de la ecuación funcional}

La derivación de $g'(0) = -\log Q$ a través de la ecuación funcional infinita es un triunfo del análisis de secuencias automáticas. Hemos logrado lo siguiente:

\begin{itemize}
    \item \textbf{Sustituto del Producto de Euler:} La ecuación funcional infinita juega exactamente el mismo papel que el Producto de Euler para la función Zeta de Riemann. Mientras que el Producto de Euler explota la multiplicatividad para factorizar $\zeta(s)$, la ecuación funcional explota la 2-automaticidad para relacionar $g(s)$ consigo misma en argumentos desplazados.
    \item \textbf{Regularidad en el Origen:} Hemos probado que $g(s)$ es una función entera, lo que garantiza que $\zeta_O(s)$ y $\zeta_E(s)$ no tienen singularidades patológicas en $s=0$. La fórmula maestra $\prod \lambda_i = \exp(-\zeta_\Lambda'(0))$ es perfectamente aplicable.
    \item \textbf{Reducción a Constantes Exóticas:} El problema analítico de calcular el producto de todos los números odiosos (o malvados) se ha reducido completamente a evaluar la constante $Q$. 
\end{itemize}

La ecuación funcional infinita nos ha llevado hasta la puerta de la regularización exótica. Con $g(0) = -1$ y $g'(0) = -\log Q$ en mano, estamos en posición de calcular explícitamente $\zeta_O'(0)$ y $\zeta_E'(0)$, lo que nos dará los productos Zeta-regularizados de los números odiosos y malvados. La evaluación final de $Q$ y la aparición de constantes fundamentales como la constante de Euler-Mascheroni $\gamma$ y la constante de Flajolet-Martin $\phi$ será el tema central de la siguiente sección.


\section{Cálculo de productos regulados exóticos}

Habiendo establecido la continuación analítica de la función zeta de Thue-Morse $g(s)$ y su derivada en el origen mediante la ecuación funcional infinita, poseemos todas las herramientas necesarias para abordar el cálculo explícito de los productos Zeta-regularizados de los números odiosos y malvados. Este cálculo no solo arroja valores numéricos sorprendentes, sino que revela una conexión profunda entre la combinatoria de las secuencias automáticas y las constantes fundamentales del análisis clásico, como la constante de Euler-Mascheroni y la constante de Flajolet-Martin.

\subsection{Evaluación del producto de todos los números odiosos}

El objetivo central es evaluar el producto Zeta-regularizado de todos los números odiosos, el cual se define formalmente mediante la fórmula maestra como:

$$ \prod_{n \in O}^{\text{reg}} n := \exp(-\zeta_O'(0)) $$

En la sección anterior, demostramos que la función zeta de los números odiosos se descompone como $\zeta_O(s) = \frac{1}{2}\zeta(s) - \frac{1}{2}g(s)$. Derivando esta expresión respecto a $s$ y evaluando en $s=0$, obtenemos la relación fundamental:

$$ \zeta_O'(0) = \frac{1}{2}\zeta'(0) - \frac{1}{2}g'(0) $$

El primer término es un resultado clásico y rigurosamente establecido de la función Zeta de Riemann: $\zeta'(0) = -\frac{1}{2}\log(2\pi)$. El desafío analítico reside por completo en evaluar $g'(0)$. 

Para calcular $g'(0)$, partimos de la ecuación funcional infinita derivada previamente para la función zeta de Thue-Morse:

$$ g(s) = -1 + \sum_{k=1}^{\infty} (-1)^{k+1} \binom{s+k-1}{k} \frac{g(s+k)}{2^{s+k}} $$

Derivamos ambos lados respecto a $s$. Dado que el coeficiente binomial $\binom{s+k-1}{k} = \frac{s(s+1)\dots(s+k-1)}{k!}$ tiene una raíz en $s=0$, su valor en $s=0$ es cero para todo $k \ge 1$. Al aplicar la regla del producto, el único término que sobrevive al evaluar en $s=0$ es aquel donde la derivada actúa sobre el factor $s$ del numerador del coeficiente binomial. La derivada de este coeficiente evaluada en $s=0$ es exactamente $\frac{1}{k}$. Por lo tanto, la derivada de la serie en el origen se simplifica drásticamente a:

$$ g'(0) = \sum_{k=1}^{\infty} (-1)^{k+1} \frac{1}{k} \frac{g(k)}{2^k} $$

Para llevar esta expresión a una forma cerrada, expandimos $g(k)$ usando su definición original como serie de Dirichlet, $g(k) = \sum_{n=1}^{\infty} \varepsilon_n n^{-k}$, y sustituimos en la ecuación:

$$ g'(0) = \sum_{k=1}^{\infty} \frac{(-1)^{k+1}}{k 2^k} \sum_{n=1}^{\infty} \frac{\varepsilon_n}{n^k} $$

Dado que las series involucradas convergen absolutamente, podemos intercambiar el orden de las sumatorias sin alterar el resultado:

$$ g'(0) = \sum_{n=1}^{\infty} \varepsilon_n \sum_{k=1}^{\infty} \frac{(-1)^{k+1}}{k} \left(\frac{1}{2n}\right)^k $$

Reconocemos inmediatamente que la suma interna es la expansión en serie de Taylor de la función logaritmo natural, $\log(1+x) = \sum_{k=1}^{\infty} \frac{(-1)^{k+1}}{k} x^k$, evaluada en $x = \frac{1}{2n}$. Esto nos permite colapsar la suma infinita en una expresión logarítmica exacta:

$$ g'(0) = \sum_{n=1}^{\infty} \varepsilon_n \log\left(1 + \frac{1}{2n}\right) = \sum_{n=1}^{\infty} \varepsilon_n \log\left(\frac{2n+1}{2n}\right) $$

Utilizando las propiedades de los logaritmos, esta suma se puede reescribir como el logaritmo de un producto infinito. Definimos la constante exótica $Q$ como el producto infinito inverso:

$$ Q := \prod_{n=1}^{\infty} \left(\frac{2n}{2n+1}\right)^{\varepsilon_n} $$

Por lo tanto, hemos demostrado rigurosamente que $g'(0) = -\log Q$. 

Sustituyendo este resultado y el valor de $\zeta'(0)$ en nuestra ecuación para $\zeta_O'(0)$, obtenemos:

$$ \zeta_O'(0) = \frac{1}{2}\left(-\frac{1}{2}\log(2\pi)\right) - \frac{1}{2}(-\log Q) = -\frac{1}{4}\log(2\pi) + \frac{1}{2}\log Q $$

Finalmente, aplicamos la fórmula maestra para obtener el producto regularizado de los números odiosos:

$$ \prod_{n \in O}^{\text{reg}} n = \exp\left( \frac{1}{4}\log(2\pi) - \frac{1}{2}\log Q \right) = (2\pi)^{1/4} Q^{-1/2} $$

De manera completamente análoga, para los números malvados, cuya función zeta es $\zeta_E(s) = \frac{1}{2}\zeta(s) + \frac{1}{2}g(s)$, la derivada en el origen es $\zeta_E'(0) = -\frac{1}{4}\log(2\pi) - \frac{1}{2}\log Q$, lo que arroja su producto regularizado:

$$ \prod_{n \in E}^{\text{reg}} n = (2\pi)^{1/4} Q^{1/2} $$

\subsection{Aparición de la constante de Euler-Mascheroni y la constante de Flajolet-Martin}

El resultado anterior es analíticamente satisfactorio, pero deja abierta la pregunta sobre la naturaleza de la constante $Q$. ¿Existe una forma cerrada para este producto infinito ponderado por la secuencia de Thue-Morse? La respuesta nos conecta con dos de las constantes más fascinantes de la teoría analítica de números y el análisis combinatorio.

La constante $Q$ está íntimamente ligada a la \textbf{constante de Flajolet-Martin}, denotada por $\phi$. Esta constante surge originalmente en el análisis de algoritmos probabilísticos de conteo en bases de datos (específicamente en el algoritmo de conteo de elementos distintos de Flajolet-Martin), pero su estructura analítica la vincula directamente con la secuencia de Thue-Morse. La constante de Flajolet-Martin se define mediante un producto infinito similar a $Q$, pero con una estructura de bloques de 4:

$$ \phi := 2^{-1/2} e^\gamma \prod_{n=1}^{\infty} \left( \frac{(4n+1)(4n+2)}{4n(4n+3)} \right)^{\varepsilon_n} \approx 0.77351\dots $$

donde $\gamma \approx 0.57721\dots$ es la célebre \textbf{constante de Euler-Mascheroni}. 

A través de manipulaciones algebraicas de los productos infinitos y utilizando las propiedades de la secuencia de Thue-Morse, se puede demostrar la siguiente relación exacta entre $Q$, $\gamma$ y $\phi$:

$$ \phi = \frac{2^{-1/2} e^\gamma}{Q} \implies Q = \frac{2^{-1/2} e^\gamma}{\phi} $$

Esta identidad es el puente que permite expresar los productos de los números odiosos y malvados puramente en términos de constantes matemáticas fundamentales. Sustituyendo $Q$ en las fórmulas derivadas previamente, obtenemos los resultados finales y cerrados:

\begin{itemize}
    \item \textbf{Producto de los números odiosos:}
    $$ \prod_{n \in O}^{\text{reg}} n = (2\pi)^{1/4} \left( \frac{2^{-1/2} e^\gamma}{\phi} \right)^{-1/2} = \pi^{1/4} \sqrt{2} \, \phi^{1/2} e^{-\gamma/2} $$
    
    \item \textbf{Producto de los números malvados:}
    $$ \prod_{n \in E}^{\text{reg}} n = (2\pi)^{1/4} \left( \frac{2^{-1/2} e^\gamma}{\phi} \right)^{1/2} = \pi^{1/4} \frac{e^{\gamma/2}}{\sqrt{2} \, \phi^{1/2}} $$
\end{itemize}

\subsubsection*{La prueba definitiva de consistencia: La Fórmula de Lerch}

La belleza de la Zeta-Regularización radica en su coherencia algebraica interna. Dado que los números odiosos $O$ y los números malvados $E$ forman una partición exacta de los números naturales $\mathbb{N}$ (es decir, $O \sqcup E = \mathbb{N}$), la propiedad de multiplicatividad por particiones de conjuntos (establecida en la Clase 1) exige que el producto de todos los números naturales sea igual al producto de los odiosos multiplicado por el producto de los malvados. 

Verifiquemos si nuestros resultados exóticos satisfacen esta exigencia algebraica:

$$ \prod_{n \in O}^{\text{reg}} n \times \prod_{n \in E}^{\text{reg}} n = \left( (2\pi)^{1/4} Q^{-1/2} \right) \times \left( (2\pi)^{1/4} Q^{1/2} \right) $$

Al multiplicar ambas expresiones, las dependencias de la constante exótica $Q$ (y por ende, de $\phi$ y $\gamma$) se cancelan mágicamente:

$$ \prod_{n \in \mathbb{N}}^{\text{reg}} n = (2\pi)^{1/4 + 1/4} Q^{-1/2 + 1/2} = (2\pi)^{1/2} Q^0 = \sqrt{2\pi} $$

Hemos recuperado exactamente la Fórmula de Lerch clásica para el producto de todos los números naturales. Esta cancelación no es una coincidencia; es la manifestación de que la Zeta-Regularización preserva rigurosamente la estructura aditiva y multiplicativa de los conjuntos subyacentes, incluso cuando los productos individuales involucran constantes trascendentales complejas como $\gamma$ y $\phi$.

\subsubsection*{Significado analítico}

La aparición de la constante de Euler-Mascheroni $\gamma$ y la constante de Flajolet-Martin $\phi$ en el producto de los números odiosos es un fenómeno analítico profundo. En la regularización de los números naturales, el producto $\sqrt{2\pi}$ surge de la función Gamma y la ecuación funcional de Riemann. Sin embargo, al restringir el producto a un subconjunto definido por la paridad de los dígitos binarios, la estructura multiplicativa clásica se rompe y es reemplazada por la autosimilitud fractal de la secuencia de Thue-Morse. 

La ecuación funcional infinita actúa como un "Producto de Euler" para secuencias no multiplicativas, y al evaluarla en el origen, la suma de las series logarítmicas resultantes cristaliza en la constante $Q$. La relación $Q = 2^{-1/2} e^\gamma / \phi$ demuestra que la constante de Flajolet-Martin no es solo una curiosidad de la informática teórica, sino que es la constante analítica natural que gobierna la regularización de productos sobre conjuntos definidos por propiedades de base 2. La Zeta-Regularización, por lo tanto, unifica la combinatoria digital, el análisis complejo y la teoría de números en un solo marco coherente.

\section{Generalizaciones y secuencias desplazadas}

El estudio de los números odiosos y malvados nos ha proporcionado un marco robusto para regularizar productos infinitos sobre conjuntos definidos por propiedades combinatorias de sus dígitos binarios. Sin embargo, la verdadera potencia y flexibilidad de la Zeta-Regularización se pone a prueba cuando sometemos a estos conjuntos a transformaciones algebraicas elementales, como los desplazamientos aditivos. En esta sección, analizaremos rigurosamente los conjuntos desplazados $OS = \{n+1 \mid n \in O\}$ y $ES = \{n+1 \mid n \in E\}$. Este ejercicio no solo generaliza los resultados anteriores, sino que introduce de manera natural una función analítica fundamental: la función $f(s)$ de Hurwitz generalizada, la cual codifica la secuencia de Thue-Morse desplazada.

\subsection{Definición de los conjuntos desplazados y partición de los enteros}

Consideremos los conjuntos de los números odiosos $O$ y malvados $E$, los cuales particionan los enteros positivos $\mathbb{N} = \{1, 2, 3, \dots\}$. Al aplicar un desplazamiento aditivo de $+1$ a cada elemento, obtenemos los nuevos conjuntos:

\begin{itemize}
    \item $OS = \{n+1 \mid n \in O\} = \{2, 3, 5, 8, 9, 12, 14, 15, 17, 20, 22, 23, 26, 27, 29, 32, \dots\}$
    \item $ES = \{n+1 \mid n \in E\} = \{4, 6, 7, 10, 11, 13, 16, 18, 19, 21, 24, 25, 28, 30, 31, \dots\}$
\end{itemize}

Es inmediato verificar que $OS$ y $ES$ forman una partición exacta de los enteros mayores o iguales a 2. Es decir, $OS \sqcup ES = \{2, 3, 4, \dots\} = \mathbb{N} \setminus \{1\}$. 

El objetivo analítico es calcular los productos Zeta-regularizados de estos nuevos conjuntos:
$$ \prod_{n \in OS}^{\text{reg}} n := \exp(-\zeta'_{OS}(0)), \qquad \prod_{n \in ES}^{\text{reg}} n := \exp(-\zeta'_{ES}(0)) $$

Para lograrlo, debemos construir las funciones zeta asociadas $\zeta_{OS}(s)$ y $\zeta_{ES}(s)$ y demostrar cómo se relacionan con la función Zeta de Riemann y la nueva función $f(s)$.

\subsection{Construcción de $\zeta_{OS}(s)$ y la función $f(s)$}

Por definición, la función zeta asociada al conjunto desplazado $OS$ es:
$$ \zeta_{OS}(s) = \sum_{n \in O} \frac{1}{(n+1)^s} = \sum_{n=1}^{\infty} \frac{a(n)}{(n+1)^s} $$
donde $a(n)$ es la función característica de los números odiosos. Utilizando la relación $a(n) = \frac{1 - \varepsilon_n}{2}$, donde $\varepsilon_n = (-1)^{a(n)}$ es la secuencia de Thue-Morse, podemos reescribir la serie como:
$$ \zeta_{OS}(s) = \frac{1}{2} \sum_{n=1}^{\infty} \frac{1}{(n+1)^s} - \frac{1}{2} \sum_{n=1}^{\infty} \frac{\varepsilon_n}{(n+1)^s} $$

Analicemos cada sumatoria por separado. La primera sumatoria es simplemente la función Zeta de Riemann a la que se le ha extraído el primer término ($n=1$):
$$ \sum_{n=1}^{\infty} \frac{1}{(n+1)^s} = \sum_{m=2}^{\infty} \frac{1}{m^s} = \zeta(s) - 1 $$

Para la segunda sumatoria, introducimos la \textbf{función $f(s)$ de Hurwitz generalizada para la secuencia de Thue-Morse}, la cual se define incluyendo el término $n=0$ (recordando que $\varepsilon_0 = 1$):
$$ f(s) := \sum_{n=0}^{\infty} \frac{\varepsilon_n}{(n+1)^s} = 1 + \sum_{n=1}^{\infty} \frac{\varepsilon_n}{(n+1)^s} $$
De aquí se deduce que $\sum_{n=1}^{\infty} \frac{\varepsilon_n}{(n+1)^s} = f(s) - 1$.

Sustituyendo ambas expresiones en la ecuación de $\zeta_{OS}(s)$, obtenemos una cancelación exacta de los términos constantes:
$$ \zeta_{OS}(s) = \frac{1}{2} (\zeta(s) - 1) - \frac{1}{2} (f(s) - 1) = \frac{1}{2} \zeta(s) - \frac{1}{2} f(s) $$

De manera completamente análoga, para el conjunto desplazado de los números malvados $ES$, cuya función característica es $1 - a(n) = \frac{1 + \varepsilon_n}{2}$, la función zeta es:
$$ \zeta_{ES}(s) = \sum_{n=1}^{\infty} \frac{1 - a(n)}{(n+1)^s} = \frac{1}{2} (\zeta(s) - 1) + \frac{1}{2} (f(s) - 1) = \frac{1}{2} \zeta(s) + \frac{1}{2} f(s) - 1 $$

\subsection{Continuación analítica y evaluación de $f'(0)$}

La función $f(s)$ es una serie de Dirichlet cuyos coeficientes son la secuencia de Thue-Morse desplazada. Al igual que la función $g(s)$ estudiada en la sección anterior, la 2-automaticidad de la secuencia de Thue-Morse garantiza que $f(s)$ admite una continuación analítica a todo el plano complejo $\mathbb{C}$ mediante una ecuación funcional infinita adaptada al desplazamiento.

Para calcular el producto regularizado, necesitamos evaluar la derivada de $f(s)$ en el origen, $f'(0)$. Por definición:
$$ f'(0) = -\sum_{n=0}^{\infty} \varepsilon_n \log(n+1) $$

A través del análisis de la ecuación funcional infinita para la secuencia desplazada (desarrollado en la literatura de series de Dirichlet automáticas por Allouche y Cohen), se demuestra rigurosamente que el desplazamiento aditivo introduce un factor de escala asintótico que modifica el valor de la derivada en el origen. El resultado exacto es:
$$ f'(0) = \frac{\log 2}{2} $$

Este valor es fundamental: refleja la asimetría combinatoria introducida al desplazar la secuencia. Mientras que la función $g(s)$ (sin desplazar) satisfacía $g(0) = -1$ y $g'(0) = -\log Q$, la función desplazada $f(s)$ absorbe las fluctuaciones del primer término, resultando en una derivada puramente logarítmica y racional.

\subsection{Cálculo de los productos regularizados exóticos}

Con las descomposiciones de $\zeta_{OS}(s)$ y $\zeta_{ES}(s)$, y el valor de $f'(0)$, estamos en posición de calcular los productos regularizados.

\subsubsection*{Producto de los números odiosos desplazados ($OS$)}

Derivamos la función $\zeta_{OS}(s) = \frac{1}{2}\zeta(s) - \frac{1}{2}f(s)$ y evaluamos en $s=0$:
$$ \zeta'_{OS}(0) = \frac{1}{2}\zeta'(0) - \frac{1}{2}f'(0) $$

Sustituimos los valores clásicos $\zeta'(0) = -\frac{1}{2}\log(2\pi)$ y $f'(0) = \frac{\log 2}{2}$:
$$ \zeta'_{OS}(0) = \frac{1}{2}\left(-\frac{1}{2}\log(2\pi)\right) - \frac{1}{2}\left(\frac{\log 2}{2}\right) = -\frac{1}{4}\log(2\pi) - \frac{1}{4}\log 2 $$

Utilizando las propiedades de los logaritmos, combinamos los términos:
$$ \zeta'_{OS}(0) = -\frac{1}{4}(\log(2\pi) + \log 2) = -\frac{1}{4}\log(4\pi) $$

Aplicando la fórmula maestra, el producto regularizado de $OS$ es:
$$ \prod_{n \in OS}^{\text{reg}} n = \exp(-\zeta'_{OS}(0)) = \exp\left(\frac{1}{4}\log(4\pi)\right) = (4\pi)^{1/4} = \sqrt{2}\pi^{1/4} $$

\subsubsection*{Producto de los números malvados desplazados ($ES$)}

Para $ES$, derivamos $\zeta_{ES}(s) = \frac{1}{2}\zeta(s) + \frac{1}{2}f(s) - 1$. Notamos que la derivada de la constante $-1$ es cero, por lo que:
$$ \zeta'_{ES}(0) = \frac{1}{2}\zeta'(0) + \frac{1}{2}f'(0) $$

Sustituyendo los mismos valores:
$$ \zeta'_{ES}(0) = -\frac{1}{4}\log(2\pi) + \frac{1}{4}\log 2 = -\frac{1}{4}(\log(2\pi) - \log 2) = -\frac{1}{4}\log(\pi) $$

El producto regularizado de $ES$ es entonces:
$$ \prod_{n \in ES}^{\text{reg}} n = \exp(-\zeta'_{ES}(0)) = \exp\left(\frac{1}{4}\log(\pi)\right) = \pi^{1/4} $$

\subsection{Verificación de consistencia algebraica: La Fórmula de Lerch}

Como en toda la teoría de la Zeta-Regularización, los resultados deben ser coherentes con las propiedades algebraicas de los conjuntos subyacentes. Dado que $OS$ y $ES$ son disjuntos y su unión es $\{2, 3, 4, \dots\}$, la propiedad de multiplicatividad por particiones de conjuntos exige que:
$$ \prod_{n \in OS}^{\text{reg}} n \times \prod_{n \in ES}^{\text{reg}} n = \prod_{n=2}^{\infty} n $$

Verifiquemos si nuestros resultados exóticos satisfacen esta condición:
$$ \left( \sqrt{2}\pi^{1/4} \right) \times \left( \pi^{1/4} \right) = \sqrt{2} \pi^{1/2} = \sqrt{2\pi} $$

Por otro lado, el producto regularizado de todos los enteros a partir de 2 se obtiene directamente de la Fórmula de Lerch clásica, dividiendo el producto de todos los naturales entre el primer elemento (que es 1):
$$ \prod_{n=2}^{\infty} n = \frac{\prod_{n=1}^{\infty} n}{1} = \frac{\sqrt{2\pi}}{1} = \sqrt{2\pi} $$

La coincidencia es exacta. Esta verificación no es trivial; confirma que la introducción de la función $f(s)$ y el desplazamiento aditivo no rompen la estructura analítica subyacente, sino que la preservan de manera rigurosa.

\subsection{Significado analítico de los desplazamientos}

El análisis de los conjuntos $OS$ y $ES$ revela una propiedad profunda de la Zeta-Regularización: la sensibilidad a la estructura fina de la secuencia. 

\begin{enumerate}[a)]
    \item \textbf{El papel de $f(s)$:} La función $f(s)$ actúa como un "potencial de desplazamiento". Mientras que la función $g(s)$ codifica la oscilación pura de la secuencia de Thue-Morse, $f(s)$ absorbe el impacto de alterar el índice de la serie. El valor $f'(0) = \frac{\log 2}{2}$ es la huella analítica de este desplazamiento.
    \item \textbf{Modificación de las constantes exóticas:} A diferencia de los conjuntos originales $O$ y $E$, cuyos productos dependían de la constante de Flajolet-Martin $\phi$ (y por ende de la constante $Q$), los productos de los conjuntos desplazados $OS$ y $ES$ son notablemente más "simples" en su forma cerrada ($\sqrt{2}\pi^{1/4}$ y $\pi^{1/4}$). Esto sugiere que el desplazamiento aditivo, al alinear la secuencia de Thue-Morse con los polos de la función Zeta de Hurwitz, cancela las contribuciones más complejas de la constante $Q$, dejando solo raíces de números enteros y $\pi$.
    \item \textbf{Robustez del formalismo:} El hecho de que podamos calcular productos sobre progresiones aritméticas, primos, secuencias automáticas y sus desplazamientos, manteniendo en todos los casos la coherencia con la Fórmula de Lerch, demuestra que la Zeta-Regularización es un marco unificado y autosuficiente. No requiere ajustes ad hoc; la continuación analítica y las ecuaciones funcionales se encargan de extraer el valor finito correcto, respetando siempre la aritmética subyacente del conjunto regularizado.
\end{enumerate}

Con este análisis de secuencias desplazadas, cerramos el estudio de las aplicaciones combinatorias y automáticas de la Zeta-Regularización, habiendo demostrado su capacidad para asignar valores finitos y algebraicamente consistentes a las estructuras más irregulares y fractales de los números enteros.

